\documentclass[a4paper,12pt]{article}
\usepackage[T2A]{fontenc}
\usepackage[utf8]{inputenc}
\usepackage[russian]{babel}
\usepackage{amsmath,amssymb,mathtools}
\usepackage{cmbright}
\renewcommand{\familydefault}{\sfdefault}
\usepackage{icomma}
\usepackage{geometry}
\geometry{left=19mm,right=19mm,top=22mm,bottom=18mm}
\setlength{\headheight}{25pt}
\usepackage{microtype}
\usepackage{tikz}
\usetikzlibrary{calc,arrows.meta,positioning,shapes.geometric}
\usepackage{tabularx,booktabs,longtable}
\usepackage{enumitem}
\usepackage{hyperref}
\usepackage{parskip}
\usepackage{fancyhdr}
\usepackage{setspace}
\usepackage{xcolor}
\usepackage{array}
\usepackage{needspace}

\definecolor{accent}{HTML}{176B6B}
\definecolor{darkaccent}{HTML}{0B4545}
\definecolor{textgray}{HTML}{2F3437}
\definecolor{softgray}{HTML}{6B7378}
\definecolor{linegray}{HTML}{D8DEDE}

\hypersetup{colorlinks=true,linkcolor=darkaccent,urlcolor=darkaccent}
\setlength{\parindent}{0pt}
\onehalfspacing
\setlist[itemize]{leftmargin=6mm,itemsep=1.5mm,topsep=1mm}
\setlist[enumerate]{leftmargin=8mm,itemsep=2.2mm,topsep=1mm}
\renewcommand{\arraystretch}{1.25}

\pagestyle{fancy}
\fancyhf{}
\fancyhead[L]{\small\color{softgray}Тригонометрические преобразования}
\fancyhead[R]{\small\color{softgray}05.10.26}
\fancyfoot[C]{\small\color{softgray}\thepage}
\renewcommand{\headrulewidth}{0.3pt}
\renewcommand{\headrule}{\hbox to\headwidth{\color{linegray}\leaders\hrule height \headrulewidth\hfill}}

\newcommand{\sectionline}{\vspace{-1mm}\noindent\color{accent}\rule{\linewidth}{0.7pt}\color{textgray}\vspace{2mm}}
\newcommand{\key}[1]{\textbf{\color{darkaccent}#1}}
\newcommand{\formula}[1]{\[\displaystyle #1\]}

\begin{document}
\color{textgray}

% ---------- TITLE ----------
\thispagestyle{empty}
\begin{center}
\vspace*{22mm}
{\Large\color{softgray}Профильная математика · ЕГЭ\par}
\vspace{9mm}
{\fontsize{29}{34}\selectfont\bfseries\color{darkaccent}Чек-лист\par}
\vspace{4mm}
{\fontsize{20}{24}\selectfont\bfseries Тригонометрические преобразования\par}
\vspace{3mm}
{\large Формулы, знаки и рациональный выбор метода\par}

\vspace{18mm}
\begin{tabular}{@{}rl@{}}
\color{softgray}Ученик: & \textbf{Ярослав Гаврилов}\\[2mm]
\color{softgray}Дата: & \textbf{5 октября 2026}\\[2mm]
\color{softgray}Домашняя работа: & \textbf{формулы: 5/день · упражнения: 3/день}
\end{tabular}

\vfill
{\large\bfseries\color{darkaccent}Лёвин Артём Александрович\par}
\vspace{1mm}
{\normalsize эксклюзивно для Ярослава Гаврилова\par}
\vspace{12mm}

\begin{tikzpicture}[x=1cm,y=1cm]
  \draw[accent,line width=1.2pt]
    (-7,0) .. controls (-4.5,0.7) and (-2.2,-0.55) .. (0,0)
    .. controls (2.2,0.55) and (4.5,-0.7) .. (7,0);
  \draw[linegray,line width=0.5pt]
    (-6.2,-0.35) .. controls (-3.8,0.15) and (-1.8,-0.75) .. (0,-0.25)
    .. controls (1.8,0.25) and (3.8,-0.55) .. (6.2,-0.15);
\end{tikzpicture}
\end{center}
\clearpage

% ---------- CORE ----------
\section*{1. Главная идея занятия}
\sectionline

В вычислительной тригонометрии скорость появляется тогда, когда Вы \key{сначала приводите выражение к удобной структуре}, а уже затем считаете. Главный вопрос перед преобразованием:

\begin{center}
\large\bfseries\color{darkaccent}
«Какая формула прямо сейчас делает выражение проще и использует данные задачи?»
\end{center}

Рабочая логика:
\begin{enumerate}
  \item сравните углы в выражении;
  \item если они отличаются в два раза, проверьте формулы двойного угла;
  \item если присутствует сдвиг на $\frac{\pi}{2}$, $\pi$, $\frac{3\pi}{2}$ и т.д., проверьте формулы приведения;
  \item выбирайте ту форму тождества, где уже есть известная величина;
  \item после преобразования оцените: стало ли выражение проще;
  \item при извлечении корня и при формулах приведения обязательно контролируйте знак.
\end{enumerate}

\section*{2. Формулы, которые должны вспоминаться мгновенно}
\sectionline

\begin{tabularx}{\linewidth}{@{}>{\bfseries\color{darkaccent}}p{0.30\linewidth}X@{}}
\toprule
Основное тождество & $\sin^2\alpha+\cos^2\alpha=1$ \\
\midrule
Тангенс и котангенс & $\displaystyle \tg\alpha=\frac{\sin\alpha}{\cos\alpha}$, $\cos\alpha\ne0$; \quad $\displaystyle \ctg\alpha=\frac{\cos\alpha}{\sin\alpha}=\frac{1}{\tg\alpha}$, где выражения определены. \\
\midrule
Синус двойного угла & $\sin 2\alpha=2\sin\alpha\cos\alpha$ \\
\midrule
Косинус двойного угла & $\cos 2\alpha=\cos^2\alpha-\sin^2\alpha=2\cos^2\alpha-1=1-2\sin^2\alpha$ \\
\bottomrule
\end{tabularx}

\vspace{4mm}
\key{Как выбирать форму для $\cos 2\alpha$.} Если дан $\sin\alpha$, чаще всего выгодна формула $1-2\sin^2\alpha$. Если дан $\cos\alpha$, выгодна $2\cos^2\alpha-1$. Это избавляет от лишнего вычисления второй функции.

\section*{3. Знаки по четвертям}
\sectionline

\begin{center}
\begin{tabular}{@{}c c c c@{}}
\toprule
Четверть & $\sin$ & $\cos$ & $\tg$ \\
\midrule
I & $+$ & $+$ & $+$ \\
II & $+$ & $-$ & $-$ \\
III & $-$ & $-$ & $+$ \\
IV & $-$ & $+$ & $-$ \\
\bottomrule
\end{tabular}
\end{center}

Если из $\sin^2\alpha+\cos^2\alpha=1$ получилось
\[
\cos^2\alpha=c,
\]
то сначала пишите
\[
\cos\alpha=\pm\sqrt{c},
\]
и только после этого выбирайте знак по указанному промежутку.

\clearpage

% ---------- METHODS & EXAMPLES ----------
\section*{4. Формулы приведения: два вопроса}
\sectionline

Для выражений вида $f\!\left(k\frac{\pi}{2}\pm\alpha\right)$ действуйте в два шага.

\begin{enumerate}
  \item \key{Меняется ли функция?} При нечётном числе четвертей ($\frac{\pi}{2}$, $\frac{3\pi}{2}$, $\frac{5\pi}{2}$, \dots) функция меняется на кофункцию: $\sin\leftrightarrow\cos$, $\tg\leftrightarrow\ctg$.
  \item \key{Какой знак?} Определяйте его по знаку \emph{исходной} функции в той четверти, куда попадает угол.
\end{enumerate}

Примеры ключевых преобразований:
\[
\sin\left(\frac{7\pi}{2}-\alpha\right)=-\cos\alpha,
\qquad
\tg\left(\alpha+\frac{5\pi}{2}\right)=-\ctg\alpha.
\]

\section*{5. Четыре опорных приёма}
\sectionline

\subsection*{Приём 1. Используйте форму двойного угла под данные}
Если $\sin\alpha=-0{,}2$, то для $\cos 2\alpha$ нет смысла сначала искать $\cos\alpha$:
\[
\cos 2\alpha=1-2\sin^2\alpha
=1-2\cdot(-0{,}2)^2
=1-0{,}08=0{,}92.
\]

\key{Важно:} при подстановке отрицательного значения используйте скобки: $(-0{,}2)^2$, а не $-0{,}2^2$.

\subsection*{Приём 2. Превратите число в тригонометрию}
Пусть
\[
5\sin^2\alpha+13\cos^2\alpha=6.
\]
Нужно найти $\tg^2\alpha$. Представим $6$ как $6(\sin^2\alpha+\cos^2\alpha)$:
\begin{align*}
5\sin^2\alpha+13\cos^2\alpha
 &=6\sin^2\alpha+6\cos^2\alpha,\\
7\cos^2\alpha&=\sin^2\alpha.
\end{align*}
Здесь $\cos\alpha\ne0$ (иначе исходное равенство дало бы $5=6$), поэтому можно разделить на $\cos^2\alpha$:
\[
\tg^2\alpha=7.
\]

\subsection*{Приём 3. Если дан тангенс, свяжите синус и косинус}
Если $\tg\alpha=-2{,}5$, то
\[
\frac{\sin\alpha}{\cos\alpha}=-2{,}5
\quad\Longrightarrow\quad
\sin\alpha=-2{,}5\cos\alpha.
\]
Для выражения
\[
\frac{10\cos\alpha+4\sin\alpha+15}
{2\sin\alpha+5\cos\alpha+3}
\]
эта подстановка даёт взаимное уничтожение тригонометрических слагаемых:
\[
\frac{10\cos\alpha-10\cos\alpha+15}
{-5\cos\alpha+5\cos\alpha+3}=\frac{15}{3}=5.
\]

\subsection*{Приём 4. Углы должны «разговаривать на одном языке»}
Если встречаются, например, $\sin 6\alpha$ и функции угла $3\alpha$, используйте
\[
\sin 6\alpha=2\sin 3\alpha\cos 3\alpha.
\]
После этого появляются одинаковые углы и часто становится возможным сокращение.

\clearpage

% ---------- FLOWCHART ----------
\thispagestyle{fancy}
\begin{center}
{\LARGE\bfseries\color{darkaccent}Алгоритм: как выбрать тригонометрическое преобразование}\par
\vspace{9mm}

\begin{tikzpicture}[
  node distance=7mm and 10mm,
  startstop/.style={ellipse,draw=accent,line width=0.9pt,minimum width=44mm,minimum height=10mm,align=center,fill=white},
  process/.style={rectangle,rounded corners=2mm,draw=accent,line width=0.9pt,text width=102mm,minimum height=11mm,align=center,inner sep=2.5mm,fill=white},
  side/.style={rectangle,rounded corners=2mm,draw=accent,line width=0.9pt,text width=50mm,minimum height=13mm,align=center,inner sep=2.5mm,fill=white},
  decision/.style={diamond,aspect=2.7,draw=accent,line width=0.9pt,text width=45mm,align=center,inner sep=1.5mm,fill=white},
  arr/.style={-{Stealth[length=2.6mm]},line width=0.9pt,color=darkaccent},
  lab/.style={font=\small\bfseries,color=darkaccent,fill=white,inner sep=1pt}
]
\node[startstop] (start) {Получили выражение};
\node[process,below=of start] (angles) {Сравните углы и структуру: $\alpha$, $2\alpha$, $3\alpha$; сдвиги на $k\pi/2$; наличие $\sin^2$ и $\cos^2$};
\node[decision,below=of angles] (different) {Нужно согласовать углы?};
\node[side,below left=9mm and 5mm of different] (normalize) {Да: двойной угол или формула приведения};
\node[side,below right=9mm and 5mm of different] (direct) {Нет: сопоставьте данное и искомое};
\node[process,below=18mm of different] (choose) {Выберите форму тождества, которая использует известную величину напрямую и создаёт минимум новых вычислений};
\node[process,below=of choose] (subst) {Подставляйте аккуратно: отрицательные числа — в скобках; при извлечении корня сначала пишите $\pm$};
\node[decision,below=of subst] (simpler) {Стало проще?};
\node[side,below left=9mm and 5mm of simpler] (alt) {Нет: попробуйте другую формулу};
\node[side,below right=9mm and 5mm of simpler] (check) {Да: досчитайте; проверьте знак, ОДЗ и деление};
\node[startstop,below=18mm of simpler] (finish) {Запишите ответ};

\draw[arr] (start) -- (angles);
\draw[arr] (angles) -- (different);
\draw[arr] (different) -- node[lab,pos=0.28,left] {да} (normalize);
\draw[arr] (different) -- node[lab,pos=0.28,right] {нет} (direct);
\draw[arr] (normalize) |- (choose);
\draw[arr] (direct) |- (choose);
\draw[arr] (choose) -- (subst);
\draw[arr] (subst) -- (simpler);
\draw[arr] (simpler) -- node[lab,pos=0.28,left] {нет} (alt);
\draw[arr] (simpler) -- node[lab,pos=0.28,right] {да} (check);
\draw[arr] (alt.west) -- ++(-4mm,0) |- (choose.west);
\draw[arr] (check) |- (finish);
\end{tikzpicture}
\end{center}
\clearpage

% ---------- ERRORS ----------
\section*{6. Типичные ошибки, которые нужно отсечь}
\sectionline

\begin{enumerate}
  \item \key{Потеря второго знака после извлечения корня.} Из $\cos^2\alpha=c$ следует $\cos\alpha=\pm\sqrt c$. Интервал выбирает нужный вариант позже.
  \item \key{Знак в формуле приведения определён по новой функции.} Знак проверяется по \emph{исходной} функции и четверти.
  \item \key{Отрицательное число подставлено без скобок.} Например, в квадрате: $(-0{,}4)^2=0{,}16$.
  \item \key{Выбрана формула, которая создаёт лишнюю неизвестную.} Если дан $\sin\alpha$, формула $\cos2\alpha=1-2\sin^2\alpha$ обычно короче, чем путь через $\cos\alpha$.
  \item \key{Деление на $\sin\alpha$ или $\cos\alpha$ без проверки.} Перед делением убедитесь, что знаменатель не равен нулю; часто это можно вывести прямо из исходного условия.
  \item \key{Формула применена, а выражение стало сложнее.} Это сигнал остановиться и проверить другой путь.
\end{enumerate}

\section*{7. Мини-самопроверка перед ответом}
\sectionline

\begin{itemize}
  \item Углы после преобразования стали удобнее?
  \item Я использовал данные задачи напрямую?
  \item Все отрицательные подстановки взяты в скобки?
  \item Если извлекал корень, сначала написал $\pm$?
  \item Если выбирал знак, проверил четверть?
  \item Если делил на выражение, проверил, что оно не равно нулю?
  \item После применения формулы стало действительно проще?
\end{itemize}

\vspace{4mm}
\key{Ежедневная микропрактика формул.} На ближайшие два дня задача простая: каждый день 3--5 минут воспроизводить формулы без подсказки и сразу применять их в коротких преобразованиях. Такой режим закрепляет формулы через регулярное извлечение из памяти.

\clearpage

% ---------- DAILY FORMULAS ----------
\section*{Ежедневная тренировка формул}
\sectionline

\textbf{Режим работы.} Это отдельная короткая тренировка на запоминание. Выполняйте \emph{без} подсматривания в чек-лист, затем сразу сверяйте. На каждый день --- ровно 5 заданий на формулы.

\subsection*{06.10 · формулы · 5 заданий}
\begin{enumerate}[label=\textbf{Ф\arabic*.}]
  \item Восстановите формулу: \(\sin 2\alpha=\ldots\)
  \item Восстановите три равносильные формулы для \(\cos 2\alpha\).
  \item Восстановите формулу: \(\displaystyle \tg\alpha=\frac{\ldots}{\ldots}\).
  \item Восстановите формулу: \(\displaystyle \ctg\alpha=\frac{\ldots}{\ldots}=\frac{1}{\ldots}\).
  \item Восстановите основное тригонометрическое тождество: \(\ldots+\ldots=1\).
\end{enumerate}

\subsection*{07.10 · формулы · 5 заданий}
\begin{enumerate}[label=\textbf{Ф\arabic*.}]
  \item Восстановите формулу: \(\cos 2\alpha=1-\ldots\).
  \item Восстановите формулу: \(\sin 2\alpha=\ldots\).
  \item Восстановите связь: \(\displaystyle \tg\alpha=\frac{\sin\alpha}{\ldots}\), \quad \(\displaystyle \ctg\alpha=\frac{\cos\alpha}{\ldots}\).
  \item Восстановите формулу приведения: \(\displaystyle \sin\left(\frac{7\pi}{2}-\alpha\right)=\ldots\).
  \item Восстановите формулу приведения: \(\displaystyle \tg\left(\alpha+\frac{5\pi}{2}\right)=\ldots\).
\end{enumerate}

\vspace{2mm}
\key{Цель блока.} Формула должна вспоминаться сразу, без длинного вывода. Если какая-то запись не восстановилась с первой попытки, повторите именно её в конце этого же дня.

\clearpage

% ---------- DAILY EXERCISES ----------
\section*{Ежедневные упражнения}
\sectionline

\textbf{Режим работы.} Это отдельный блок применения формул. На каждый день --- ровно 3 полноценные задачи. Здесь уже важно выбрать подходящую формулу, аккуратно подставить данные и проверить знак.

\subsection*{06.10 · упражнения · 3 задания}
\begin{enumerate}[label=\textbf{У\arabic*.}]
  \item Известно, что \(\sin\alpha=-0{,}4\). Найдите \(\cos 2\alpha\), используя формулу, которая напрямую содержит \(\sin\alpha\).

  \item Найдите значение выражения
  \[
    \frac{\sin 6x}{\sin 3x\cos 3x},
  \]
  если оно определено.

  \item Известно, что
  \[
    4\sin^2\alpha+10\cos^2\alpha=5.
  \]
  Найдите \(\tg^2\alpha\).
\end{enumerate}

\subsection*{07.10 · упражнения · 3 задания}
\begin{enumerate}[label=\textbf{У\arabic*.}]
  \item Известно, что \(\sin\alpha=\dfrac45\) и \(\alpha\in\left(\dfrac\pi2,\pi\right)\). Найдите
  \[
    \sin\left(\frac{7\pi}{2}-\alpha\right).
  \]

  \item Известно, что \(\tg\alpha=-0{,}4\). Найдите
  \[
    \tg\left(\alpha+\frac{5\pi}{2}\right).
  \]

  \item Известно, что \(\tg\alpha=-2{,}5\). Найдите значение выражения
  \[
    \frac{10\cos\alpha+4\sin\alpha+15}
         {2\sin\alpha+5\cos\alpha+3}.
  \]
\end{enumerate}

\vspace{2mm}
\key{Важно.} Формульная тренировка и упражнения --- два разных блока. Сначала 5 коротких заданий на воспроизведение формул, затем 3 задачи на их применение.

\clearpage

% ---------- ANSWERS ----------
\section*{Ответы}
\sectionline

\subsection*{Тренировка формул}
{\small
\begin{tabularx}{\linewidth}{@{}>{\bfseries\color{darkaccent}}l >{\bfseries}c X@{}}
\toprule
Дата & № & Ответ \\
\midrule
06.10 & Ф1 & \(\sin2\alpha=2\sin\alpha\cos\alpha\). \\
06.10 & Ф2 & \(\cos2\alpha=\cos^2\alpha-\sin^2\alpha=1-2\sin^2\alpha=2\cos^2\alpha-1\). \\
06.10 & Ф3 & \(\displaystyle\tg\alpha=\frac{\sin\alpha}{\cos\alpha}\). \\
06.10 & Ф4 & \(\displaystyle\ctg\alpha=\frac{\cos\alpha}{\sin\alpha}=\frac1{\tg\alpha}\), где выражения определены. \\
06.10 & Ф5 & \(\sin^2\alpha+\cos^2\alpha=1\). \\
\midrule
07.10 & Ф1 & \(\cos2\alpha=1-2\sin^2\alpha\). \\
07.10 & Ф2 & \(\sin2\alpha=2\sin\alpha\cos\alpha\). \\
07.10 & Ф3 & \(\displaystyle\tg\alpha=\frac{\sin\alpha}{\cos\alpha}\), \(\displaystyle\ctg\alpha=\frac{\cos\alpha}{\sin\alpha}\). \\
07.10 & Ф4 & \(\displaystyle\sin\left(\frac{7\pi}{2}-\alpha\right)=-\cos\alpha\). \\
07.10 & Ф5 & \(\displaystyle\tg\left(\alpha+\frac{5\pi}{2}\right)=-\ctg\alpha\). \\
\bottomrule
\end{tabularx}
}

\vspace{4mm}
\subsection*{Ежедневные упражнения}
{\small
\begin{tabularx}{\linewidth}{@{}>{\bfseries\color{darkaccent}}l >{\bfseries}c X@{}}
\toprule
Дата & № & Ответ \\
\midrule
06.10 & У1 & \(\cos2\alpha=1-2\cdot(-0{,}4)^2=0{,}68\). \\
06.10 & У2 & \(2\). \\
06.10 & У3 & \(\tg^2\alpha=5\). \\
\midrule
07.10 & У1 & \(\displaystyle\frac35\). \\
07.10 & У2 & \(\displaystyle\frac52=2{,}5\). \\
07.10 & У3 & \(5\). \\
\bottomrule
\end{tabularx}
}

\vfill
\begin{center}
\color{softgray}\small Лёвин Артём Александрович · эксклюзивно для Ярослава Гаврилова
\end{center}

\end{document}
